\documentclass{article}
\usepackage[margin=1in]{geometry}
\usepackage{mathtools, amsfonts, amsthm, xcolor, listings}

\newcommand{\R}{\mathbb{R}}
\newcommand{\N}{\mathbb{N}}
\newcommand{\I}{\mathbb{I}}
\newcommand{\Q}{\mathbb{Q}}
\newcommand{\Z}{\mathbb{Z}}

\newcommand{\set}[1]{\left\{#1\right\}}
\newcommand{\ang}[1]{\left\langle #1 \right\rangle}
\newcommand{\inv}{^{-1}}

\newcommand{\normsub}{\lhd}

\newcommand{\reflection}[5]{
    \begin{itemize}
    \item Points attempted: #1
    \item Self-assessed score: #2/#1.
    \item Time spent: #3 minutes.
    \item Difficulty: #4/10.
    \item Questions/Feedback: #5.
    \end{itemize}
}


\title{HW08}
\author{Sam Ly}

\begin{document}
\maketitle

\section*{Exercise 1}
Collaborators: Sam

\noindent
Tools and resources used
\pagebreak

\section*{Exercise 2}
\reflection{1}{1}{2}{1}{None.}
HW8 Exercise 2: One thing that I would like to see on HW 9 would be some more work on burnside's theorem. Although I "know" the theorem, I do not have a full intuition for it yet.
\pagebreak

\section*{Choice Exercise 5}
We defined a (left) group action of the group \((G,*)\) on the set \(X\) as a function \(\cdot\colon G \times X \to X\) such that for all \(x \in X\) and \(g_1, g_2 \in G\)
\begin{itemize}
    \item \(e \cdot x = x\), and
    \item \(g_1 \cdot (g_2 \cdot x) = (g_1 * g_2) \cdot x\).
\end{itemize}

\begin{enumerate}
    \item {
        [4] For each of the following subparts, give an example of a group \((G,*)\), a set \(X\), and a function \(\cdot\colon G \times X \to X\) such that
        \begin{enumerate}
            \item {
                The function is a group action (i.e. it satisfies both of the above properties.)

                We can have the group SO(3) acting on the set of 3D vectors \(\R^3\),
                where the function is a matrix-vector multiplication. The 
                identity matrix \(I \in \text{SO}(3)\) satisfies the first 
                condition, and the associativity of matrix-vector products 
                satisify the second condition.
            }
            \item {
                The function satisfies the first property, but not the second.

                We can have the group \(C_2\) and the set \(\Z\), with the function 
                being addition, such that the identity in \(C_2\) corresponds to 
                \(0\) and \(r\) corresponds to \(1\). Now, the first condition
                is always met, however the second condition breaks.
            }
            \item {
                The function satisfies the second property, but not the first.

                Now, we can have the group \(C_2\) and the set \(\Z / 2\Z\), with 
                the function 
                \[
                    c \cdot x = \begin{cases}
                        1 + x   & c = e \\
                        x       & c = r
                    \end{cases},
                \]
                for \(c \in C_2\) and \(x \in \Z/ 2\Z\).
            }
            \item {
                We can The function satisfies neither property.
                
                We can have any non-trivial group \(G\) an any set \(X\) such that 
                \(|X| > 1\), where the function always maps to some \(x \in X\).
                This function breaks both conditions.
            }
        \end{enumerate}
    }
    \item {
        [3] Prove that if \( G \) is a group that acts on a set \( X \), then 
        the set 
        \[H = \{g \in G \mid g \cdot x = x \text{ for all } x \in X\}\] is a 
        normal subgroup of \( G \).

        \begin{proof}
            To begin, we note that for \(H\) to be a normal subgroup of \(G\), 
            \(H\) itself must be a group. So, we must show that \(H\) satisfies 
            the group axioms. First, associativity is inherited. Also notice that,
            by definition, \(e \in H\).

            Then, \(H\) is closed because for some \(h, h' \in H\), 
            \(h \cdot x = x\) and \(h' \cdot x = x\), so \(h \cdot (h' \cdot x) = x\).
            This means that \((h * h') \cdot x = x\), so \(h * h' \in H\). Thus,
            \(H\) is closed.

            Now, we must prove that for all \(h \in H\), \(h\inv \in H\). First, 
            we know that \(e \in H\), because \(e \cdot x = x\). Notice that 
            \(h\inv * h = e\), so \((h\inv * h) \cdot x = x\). By definition of 
            a group action, \((h\inv * h) \cdot x = h\inv \cdot h \cdot x = h\inv \cdot x = x\).
            Thus, \(h\inv \in H\).

            Now, to show that \(H \normsub G\), we must show that for any \(h, g\), 
            \(g * h * g\inv \in H\). Let use directly apply this expresion onto 
            some \(x \in X\) to get \(g \cdot h \cdot g\inv \cdot x\). Notice that 
            \(g\inv \cdot x \in X\), and \(h \cdot x' = x'\) for any \(x' \in X\).
            So, we can effectively remove \(h\) from the expression to get 
            \(g \cdot g\inv \cdot x = (g * g\inv) \cdot x = x\). Thus, \(g* h * g\inv \in H\), 
            so \(H \normsub G\).
        \end{proof}
    }
\end{enumerate}
  
  

\end{document}
